\chapter{Cronograma}
\label{ch:cron}

El cronograma cubre 17 semanas y sigue una secuencia simple: preparación, medición, implementación, análisis, validación y cierre. Las primeras semanas se usan para preparar el entorno, validar la cadena JTAG y comprobar la medición de potencia.

Después se ejecuta el experimento preliminar con bucles dominados por categoría. Esa actividad no depende todavía del clasificador RTL, así que puede hacerse pronto.

La etapa central corresponde al diseño del clasificador, su integración con los CSR del CV32E40P y la síntesis del SoC PULPissimo modificado sobre la FPGA Nexys~A7-100T. Luego se hace la caracterización por regresión de histogramas y se construye el flujo de análisis.

Las semanas finales se reservan para validar el modelo con cargas mixtas, comparar la estimación contra la medición física de referencia y cerrar la documentación. La semana 17 queda como margen ante atrasos de síntesis o temporización.

\definecolor{ganttPrep}{RGB}{198,224,255}
\definecolor{ganttMeas}{RGB}{190,230,210}
\definecolor{ganttRTL}{RGB}{255,224,178}
\definecolor{ganttAna}{RGB}{215,205,245}
\definecolor{ganttVal}{RGB}{255,205,210}
\definecolor{ganttDoc}{RGB}{220,220,220}
\definecolor{ganttHead}{RGB}{235,235,235}

\newcommand{\prep}{\cellcolor{ganttPrep}}
\newcommand{\meas}{\cellcolor{ganttMeas}}
\newcommand{\rtl}{\cellcolor{ganttRTL}}
\newcommand{\ana}{\cellcolor{ganttAna}}
\newcommand{\val}{\cellcolor{ganttVal}}
\newcommand{\doc}{\cellcolor{ganttDoc}}
\newcommand{\legendbox}[1]{\cellcolor{#1}\hspace{1.2em}}

\begin{table}[!htbp]
\caption{Cronograma general de ejecución del proyecto.}
\label{tab:cronograma}
\centering
\hfill
\scriptsize
\begin{tabular}{@{}ll@{\hspace{1em}}ll@{\hspace{1em}}ll@{}}
\legendbox{ganttPrep} & Preparación &
\legendbox{ganttMeas} & Medición &
\legendbox{ganttRTL} & Implementación \\
\legendbox{ganttAna} & Análisis &
\legendbox{ganttVal} & Validación &
\legendbox{ganttDoc} & Documentación \\
\end{tabular}
\normalsize
\vspace{0.4em}

\begingroup
\renewcommand{\arraystretch}{1.25}
\setlength{\arrayrulewidth}{1.1pt}
\setlength{\tabcolsep}{3pt}
\resizebox{\textwidth}{!}{%
\begin{tabular}{|p{5.0cm}|c|c|c|c|c|c|c|c|c|c|c|c|c|c|c|c|c|}
\hline
\rowcolor{ganttHead}
\textbf{Actividad} &
\textbf{1} & \textbf{2} & \textbf{3} & \textbf{4} &
\textbf{5} & \textbf{6} & \textbf{7} & \textbf{8} &
\textbf{9} & \textbf{10} & \textbf{11} & \textbf{12} &
\textbf{13} & \textbf{14} & \textbf{15} & \textbf{16} &
\textbf{17} \\
\hline
Revisión técnica y preparación del entorno &
\prep & \prep &  &  &  &  &  &  &  &  &  &  &  &  &  &  &  \\
\hline
Diseño y verificación del circuito de medición &
 & \meas & \meas & \meas &  &  &  &  &  &  &  &  &  &  &  &  &  \\
\hline
Experimento preliminar con bucles dominados por categoría &
 &  &  & \meas & \meas &  &  &  &  &  &  &  &  &  &  &  &  \\
\hline
Caracterización por bucles dominados por categoría &
 &  &  &  & \meas & \meas &  &  &  &  &  &  &  &  &  &  &  \\
\hline
Diseño RTL del clasificador de instrucciones &
 &  &  &  & \rtl & \rtl & \rtl &  &  &  &  &  &  &  &  &  &  \\
\hline
Integración de CSR y firmware de prueba &
 &  &  &  &  & \rtl & \rtl & \rtl &  &  &  &  &  &  &  &  &  \\
\hline
Síntesis e implementación en FPGA &
 &  &  &  &  &  &  & \rtl & \rtl &  &  &  &  &  &  &  &  \\
\hline
Caracterización por regresión de histogramas &
 &  &  &  &  &  &  &  & \ana & \ana & \ana &  &  &  &  &  &  \\
\hline
Pipeline de análisis y visualización &
 &  &  &  &  &  &  &  &  & \ana & \ana & \ana & \ana &  &  &  &  \\
\hline
Validación con cargas mixtas &
 &  &  &  &  &  &  &  &  &  &  & \val & \val & \val &  &  &  \\
\hline
Redacción, revisión y entrega final &
 &  &  &  &  &  &  &  &  &  & \doc & \doc & \doc & \doc & \doc & \doc &  \\
\hline
\end{tabular}%
}
\endgroup
\end{table}
